% Overby Industries White Paper
% Source: converted from overby-industries-vision.md
\documentclass[12pt]{article}

% Packages
\usepackage{geometry}
\usepackage{graphicx}
\usepackage{xcolor}
\usepackage{hyperref}
\usepackage{fancyhdr}
\usepackage{titlesec}
\usepackage{tocloft}
\usepackage{amsmath}
\usepackage{amssymb}
\usepackage{longtable}

% Page setup
\geometry{margin=1in}
\setlength{\parskip}{0.8em}
\setlength{\parindent}{0em}

% Header/footer
\pagestyle{fancy}
\fancyhf{}
\fancyhead[L]{Overby Industries White Paper}
\fancyhead[R]{\thepage}

% Title formatting
\titleformat{\section}{\large\bfseries}{\thesection}{1em}{}
\titleformat{\subsection}{\normalsize\bfseries}{\thesubsection}{1em}{}

% Hyperlinks style
\hypersetup{
    colorlinks=true,
    linkcolor=blue,
    urlcolor=blue,
    citecolor=blue
}

% Title Info
\title{
    \Huge \textbf{Overby Industries} \\[0.5em]
    \Large An Ethical, Zero-Waste Roadmap to Sustainable Space Industry \\[1em]
    {\large Version 1.0 -- Draft for Community Review}
}
\author{Overby Industries R\&D Division}
\date{\today}

\begin{document}

\maketitle
\thispagestyle{empty}

\begin{abstract}
Overby Industries is committed to building the resources backbone of a solar system economy without repeating the environmental damage caused by terrestrial extraction. Through reusable shuttlecraft (Starlifter II), autonomous miner/refinery pods, and distributed In-Situ Resource Utilization (ISRU) systems, we will deliver metals, aggregates, plastics, and fuels for space construction and expansion with zero debris and zero waste. Our mission is to become the ethical steward of humanity's industrial expansion into space.
\end{abstract}

\newpage
\tableofcontents
\newpage

\section{The Problem}
Earth's industrial extraction is environmentally destructive and limited by finite resources. Current space infrastructure costs are dominated by the mass of materials launched from Earth, and space mining proposals risk repeating the same mistakes unless sustainability and reclamation are designed in from the start.

\section{The Vision}
Overby Industries will:
\begin{itemize}
    \item Extract metals from M-class asteroids (nickel, iron, PGMs).
    \item Manufacture Ultra High-Performance Concrete (UHPC) and basalt composites from silicate-rich asteroids and lunar regolith.
    \item Establish a Titan Moon Base for hydrocarbons converted into plastics, polymers, and fuels.
    \item Provide propellants (H$_2$/O$_2$/CH$_4$) cracked from asteroid ice and CO$_2$.
\end{itemize}

All operations operate under a Zero-Pollution, Zero-Waste Certification. Every residue -- dust, slag, rock -- is repurposed into shielding, bricks, and aggregates.

\section{Core Technologies}

\subsection{Starlifter II Shuttle}
A fully reusable, runway-landing spacecraft with hybrid propulsion: ion thrusters for in-space maneuvering and Air-Breathing Electric Propulsion in atmosphere, Lorentz rail accelerators for boost maneuvers, and an MHD solar-wind dynamo for power and thrust. Ionic liquids serve as propellant, electrolyte, and internal thermal shielding during reentry.

\subsection{Autonomous Miner/Refinery Pods}
Anchored drill rigs paired with expandable containment bag systems (based on NASA ARM concepts) enable zero-dust excavation. In-situ refining produces ingots, captures volatiles, and converts slag into building materials.

\subsection{Power Systems}
Magnetohydrodynamic (MHD) reactor cores, graphene supercapacitors for burst thrust, aluminum-ion and ionic-liquid batteries for safe flame-proof energy storage, and solar wind particle capture for continuous trickle power.

\subsection{ISRU Output Streams}
\begin{itemize}
    \item \textbf{Metals:} Nickel, iron, PGMs.
    \item \textbf{Aggregates:} Silicates converted into UHPC and basalt fiber composites.
    \item \textbf{Polymers:} Titan methane/ethane converted into plastics.
    \item \textbf{Fuels:} O$_2$, H$_2$, CH$_4$ for depots and propulsion.
\end{itemize}

\section{Ethical and Environmental Stewardship}
Overby Industries sets the Space Zero-Waste Standard:
\begin{itemize}
    \item \textbf{Containment-first mining:} no dust, no orbital debris.
    \item \textbf{Reclamation mode:} every miner pod compacts residual material.
    \item \textbf{Closed-loop ISRU cycles:} nothing is left behind.
    \item \textbf{Transparency:} public certification of every mission.
\end{itemize}

\begin{quote}
\textit{``Humanity's future in space will be sustainable -- or it will fail. Overby Industries ensures it will succeed.''}
\end{quote}

\section{Roadmap Timeline}

\begin{longtable}{p{3cm} p{2.2cm} p{8cm}}
\textbf{Phase} & \textbf{Years} & \textbf{Milestones} \\
\hline
\endhead
Foundations & 2030--2040 & Starlifter prototypes, Miner Pod testbeds, NEA mining demo \\
Industrial Pilot & 2040--2055 & Fleet of Starlifter II, lunar depots, first profitable asteroid resources \\
Expansion & 2055--2080 & UHPC structures in orbit, Mars staging post \\
Outer System & 2080--2100 & Titan Base, plastics and fuels economy \\
Civilization Scale & 2100+ & Full solar system closed-loop infrastructure, Overby Certification becomes global standard \\
\end{longtable}

\section{Conclusion}
This white paper establishes the founding vision that the rest of the Overby Industries ecosystem -- power and propulsion research, orbital debris reclamation services, the civic governance framework, and the public-facing platform -- is built to serve. It is shared openly to invite constructive input from the global community of thinkers, builders, and explorers this mission depends on.

\section{References}
[1] Placeholder for academic references.\\{}
[2] Placeholder for community submissions.

\end{document}
